One mitigation is carried through the design and pre-registered as $M1$: increase patch overlap $O$ by reducing the stride $S$ at fixed patch size $P$. Its geometric basis lies in the stride branch, whose candidates are in $\mathcal{F}_{S}$. $M1$ asks whether this reduction also decreases the observed recovery deficit. Model~A$'$ of \autoref{sec:four} tests the claim through the overlap slope $\delta_O$, read on the contrast $d$ that Model~A is fitted to, where a mitigation is a positive slope. The pre-registered rule has two legs: $\Pr(\delta_O>0\mid D)\geq0.95$, and the configuration level carrying the overlap term must win the leave-one-out comparison against the same level without it. Both must hold.

Part of the answer does not wait on the fit. In the analysed band the stride candidates satisfy $f=cf_s/S$ with $c<S/2$, so a geometry has $\lceil S/2\rceil-1$ of them; a smaller $S$ raises the fundamental $f_s/S$, widens the spacing and lowers that count, and halving $S$ retains only the even-indexed sites. The patch fundamental $f_s/P$ is untouched by any of this, so $\mathcal{F}_{P}$ is invariant under the mitigation by construction. Overlap is therefore a partial geometric adjustment at best: it displaces stride-derived candidates and cannot remove patch-derived ones.

\textcolor{red}{\textbf{Placeholder, to be replaced once the fits are complete.} This subsection reports the posterior of $\delta_O$ on both fitted sets with its credible interval, the probability $\Pr(\delta_O>0\mid D)$ against the first leg of the rule, and the leave-one-out comparison of the configuration level with and without the overlap term against the second leg, followed by the three-way verdict and the convergence diagnostics that gate it. It then qualifies the estimate with the three confounds this design cannot separate: the token count $N=1+\lfloor(L-P)/S\rfloor$ moves with the stride, so overlap and sequence length are not independently identified; lowering $S$ changes the mixture of stride- and patch-derived candidates along the same axis; and the truncated context tail of \autoref{sec:six} is non-zero on exactly the six geometries that identify the stride branch. None of the three could create a mitigation, but any could inflate one.}

\subsection{Further strategies}

Because the candidate set is determined by $(f_s,P,S)$ and by nothing else, two further strategies follow from the same arithmetic, and both are worth stating even though neither is tested here.

The first is to move the combs rather than to thin one of them. Changing $S$ translates the stride comb and changing $P$ translates the patch comb, each by a known amount, so a deployment that knows the spectral content carried can choose a geometry whose combs fall between the components that matter instead of on them. This is a design-time control rather than a repair: it does not make any frequency more recoverable, it relocates the frequencies at which recovery is worst, and the relocation is deterministic and computable before a single forecast is made. Its limit is that it protects only what is known in advance, so is not a mitigation against spectrum changes.

The second is redundancy across geometries: two models with different geometries are blind at different frequencies. An ensemble that selects per frequency, using the observability assessment of \autoref{sec:nine} as the switch, would leave no frequency served only by a model known to be blind there. This mitigation addresses the effect rather than relocating it, and it does so precisely because it does not rely on any single tokenisation being adequate. Its costs are real: two forecasters must be calibrated against each other, the switch itself becomes a component that can be wrong, and the claim is untested here, since establishing it would need the paired contrast of \autoref{sec:four} applied to the ensemble output rather than to a single model.

$M1$ thins one branch, geometry selection moves both, and only redundancy covers what a single tokeniser cannot represent, at the price of no longer being a single model.
